\subsection{The exact budget of the supplied sign family}
\label{sec:bf-budget-comparison}

The comparison with classical curvature inequalities can be made exact
for the classes supplied by the preceding construction.  The two
optimizations below are the convex and positive-plane budgets of
\cite{LeBrunConvex}, restricted to that specified set of classes.

Let $r\in\{2,3\}$, and suppose that the simply connected symplectic
summands $X_i$ satisfy the stable-cohomotopy hypotheses of
Section~\ref{sec:bf-geography}, including $b_+(X_i)\equiv3\pmod4$.
Assume in addition that $A_i=c_1^2(X_i)>0$ for each $i$, as in both
realized families above.  On
\[
 M=\Bigl(\#_{i=1}^rX_i\Bigr)\#k\overline{\mathbb{CP}}^{\,2}
                         \#m(S^1\times S^3),
 \qquad A=\sum_{i=1}^rA_i,
\]
write $c_i=c_1(X_i)$ and let $E_1,\ldots,E_k$ be the exceptional
generators.  The connected-sum splitting gives
\[
 c_i\cdot c_j=\delta_{ij}A_i,\qquad
 c_i\cdot E_j=0,\qquad E_i\cdot E_j=-\delta_{ij}.
\]
The actual nonvanishing theorem supplies every member of
\begin{equation}\label{eq:bf-supplied-sign-set}
 \mathcal C_0=
 \left\{\sum_{i=1}^r\epsilon_i c_i+
              \sum_{j=1}^k\eta_jE_j:
             \epsilon_i,\eta_j\in\{-1,1\}\right\}.
\end{equation}
Thus the algebraic set used here consists of actual monopole classes.

\begin{proposition}[Classical budgets of the supplied classes]
\label{prop:bf-sign-budgets}
For a maximal positive subspace $H\subset H^2(M;\mathbb R)$, let
$a_H^+$ be the orthogonal projection with respect to the intersection
form.  Define
\[
 \beta_{\mathcal C_0}^2
    =\max_{v\in\operatorname{conv}(\mathcal C_0)}v^2,
 \qquad
 \alpha_{\mathcal C_0}^2
    =\inf_H\max_{a\in\mathcal C_0}(a_H^+)^2.
\]
Then
\begin{equation}\label{eq:bf-sign-budgets}
 \beta_{\mathcal C_0}^2=\alpha_{\mathcal C_0}^2=A.
\end{equation}
\end{proposition}

\begin{proof}
The convex hull is the box of vectors
$v=\sum_i t_ic_i+\sum_j u_jE_j$ with $|t_i|,|u_j|\le1$.
Orthogonality gives
\[
 v^2=\sum_i A_it_i^2-\sum_j u_j^2\le A.
\]
The vector $\sum_i c_i$ belongs to the hull, by averaging the
exceptional signs, and has square $A$.  This proves the first equality.

For any maximal positive $H$, averaging the projected squares over
all independent signs gives
\[
 \max_{a\in\mathcal C_0}(a_H^+)^2
 \ge\sum_i(c_{i,H}^+)^2+\sum_j(E_{j,H}^+)^2
 \ge\sum_i c_i^2=A.
\]
For the reverse inequality, the orthogonal complement of the negative
space spanned by the $E_j$ contains the positive space spanned by the
$c_i$.  Extend the latter to a maximal positive subspace $H_0$ in that
orthogonal complement.  Then $c_i\in H_0$, $E_{j,H_0}^+=0$, and every
member of $\mathcal C_0$ has projected square $A$.  This proves the
second equality.  The optimization ranges over abstract positive
subspaces; it makes no assertion that $H_0$ is realized by a metric.
\end{proof}

The full set $\mathcal C(M)$ of monopole classes may be larger than
$\mathcal C_0$.  Let $B$ denote either of the analogous full-set
budgets, so that $B\ge A$.  Each is a uniform projection budget:
for every metric $g$,
\[
 \max_{a\in\mathcal C(M)}(a_g^+)^2\ge B.
\]
For the positive-plane budget this follows from its definition.  For
the convex budget, projection, the triangle inequality, and
$v^2\le(v_g^+)^2$ show that every convex combination has square at
most the largest projected square of its constituent classes.

Write $B=A(1+\delta)$ with $\delta\ge0$, and retain
$T_+=A(1-h)$ and $U=3\min\{T_+,T_-\}/A$.  The classical
Ricci-energy inequality, in the present tensor-Ricci normalization, is
\begin{equation}\label{eq:bf-complete-budget-ricci}
 R\ge8\pi^2(2p-T_+).
\end{equation}
Indeed, its deficit is precisely
$D_2=\tfrac92\{R-8\pi^2(2p-T_+)\}\ge0$.
Choosing a class with $p\ge B$ gives
\[
 e:=\frac{R}{8\pi^2A}\ge1+h+2\delta.
\]
For an Einstein metric $P=0$, the two chiral Chern--Weil identities
also give $e\le U$.  Thus the full classical budget imposes the
necessary condition
\begin{equation}\label{eq:bf-complete-budget-necessary}
 \boxed{1+h+2\delta\le U.}
\end{equation}

Proposition~\ref{prop:bf-sign-budgets} proves that optimizing over all
of the supplied signs leaves the budget equal to $A$.  It does not
bound the increment $B-A$ from additional monopole classes.  For the
linear family of Section~\ref{sec:bf-linear}, where
$U-1-h=\theta h+O(h^2)$, an increment of order $Ah$ can change this
classical comparison.  An upper estimate such as $B-A=o(Ah)$ would
remove that uncertainty at first order; it remains a separate
mathematical question about the complete class set.  The separation
theorem uses the supplied actual classes and does not require such an
upper estimate.  Its comparison is therefore with the explicitly
computed budgets of $\mathcal C_0$, without a claim that every
classical obstruction on these smooth manifolds has been exhausted.
